\documentclass[10pt,aspectratio=169,serif]{beamer}

\usepackage[fontset=none, scheme=plain]{ctex}
\IfFontExistsTF{Songti SC}{%
	\setCJKmainfont{Songti SC}[AutoFakeSlant=0.2]%
	\setCJKsansfont{Songti SC}[AutoFakeSlant=0.2]%
}{%
	\setCJKmainfont{FandolSong-Regular}[Extension=.otf,BoldFont=FandolSong-Bold,AutoFakeSlant=0.2]%
	\setCJKsansfont{FandolSong-Regular}[Extension=.otf,BoldFont=FandolSong-Bold]%
}
\usepackage{amsmath,amssymb,graphicx,tikz,booktabs,multicol,calligra}
\usepackage{pgfplots}
\pgfplotsset{compat=1.18}
\usepackage{tcolorbox}
\tcbuselibrary{skins}
\usepackage{XyygU}
\usepackage{caption}
\captionsetup[figure]{name=图,labelsep=period}

\usetikzlibrary{arrows.meta,positioning,calc}

% 公式间距
\setlength{\abovedisplayskip}{0pt plus 2pt minus 2pt}
\setlength{\belowdisplayskip}{0pt plus 2pt minus 2pt}

\title[全称量词与存在量词练习]{全称量词与存在量词（练习）}
\subtitle{1.5.1 全称量词与存在量词 \quad 1.5.2 全称量词命题和存在量词命题的否定}
\author{}
\institute{}
\date{2026-10-16}

\begin{document}

\begin{frame}
  \titlepage
\end{frame}

% ========== 1.5.1 练习 第1题 ==========
\begin{frame}[t]{1.5.1 练习 第1题：判断全称量词命题的真假}
\textbf{1.} 判断下列全称量词命题的真假：

\vspace{0.5em}
(1) 每个四边形的内角和都是 $360^\circ$；

\vspace{0.3em}
(2) 任何实数都有算术平方根；

\vspace{0.3em}
(3) $\forall x \in \{y \mid y \text{ 是无理数}\}$，$x^3$ 是无理数。
\end{frame}

\begin{frame}[t]{1.5.1 练习 第1题 解答}
\begin{tcolorbox}[colback=orange!10,colframe=orange,title=\textbf{解}]
(1) \textbf{真命题。}

\vspace{0.3em}
任意四边形都可以分成两个三角形，每个三角形内角和为 $180^\circ$，所以四边形内角和为 $360^\circ$。

\vspace{0.8em}
(2) \textbf{假命题。}

\vspace{0.3em}
负数没有算术平方根。例如 $x = -1$ 是实数，但 $-1$ 没有算术平方根。

\vspace{0.8em}
(3) \textbf{假命题。}

\vspace{0.3em}
举反例：取 $x = \sqrt[3]{2}$，则 $x$ 是无理数，但 $x^3 = 2$ 是有理数。
\end{tcolorbox}
\end{frame}

% ========== 1.5.1 练习 第2题 ==========
\begin{frame}[t]{1.5.1 练习 第2题：判断存在量词命题的真假}
\textbf{2.} 判断下列存在量词命题的真假：

\vspace{0.5em}
(1) 存在一个四边形，它的两条对角线互相垂直；

\vspace{0.3em}
(2) 至少有一个整数 $n$，使得 $n^2 + n$ 为奇数；

\vspace{0.3em}
(3) $\exists x \in \{y \mid y \text{ 是无理数}\}$，$x^2$ 是无理数。
\end{frame}

\begin{frame}[t]{1.5.1 练习 第2题 解答}
\begin{tcolorbox}[colback=orange!10,colframe=orange,title=\textbf{解}]
(1) \textbf{真命题。}

\vspace{0.3em}
菱形（或正方形）的两条对角线互相垂直。

\vspace{0.8em}
(2) \textbf{假命题。}

\vspace{0.3em}
因为 $n^2 + n = n(n+1)$ 是两个连续整数的乘积，必为偶数。对任意整数 $n$，$n$ 与 $n+1$ 中必有一个偶数，所以 $n(n+1)$ 恒为偶数，不可能为奇数。

\vspace{0.8em}
(3) \textbf{真命题。}

\vspace{0.3em}
取 $x = \sqrt[3]{2}$，则 $x$ 是无理数，$x^2 = \sqrt[3]{4}$ 也是无理数。
\end{tcolorbox}
\end{frame}

% ========== 1.5.2 练习 第1题 ==========
\begin{frame}[t]{1.5.2 练习 第1题：写全称量词命题的否定}
\textbf{1.} 写出下列命题的否定：

\vspace{0.5em}
(1) $\forall n \in \mathbf{Z}$，$n \in \mathbf{Q}$；

\vspace{0.3em}
(2) 任意奇数的平方还是奇数；

\vspace{0.3em}
(3) 每个平行四边形都是中心对称图形。
\end{frame}

\begin{frame}[t]{1.5.2 练习 第1题 解答}
\begin{tcolorbox}[colback=orange!10,colframe=orange,title=\textbf{解}]
\textbf{规则：} 全称量词命题 $\forall x \in M,\ p(x)$ 的否定为 $\exists x \in M,\ \neg p(x)$。

\vspace{0.8em}
(1) 原命题的否定：
\[
\exists n \in \mathbf{Z},\ n \notin \mathbf{Q}.
\]

\vspace{0.5em}
(2) 原命题的否定：
存在一个奇数，它的平方不是奇数（即存在一个奇数，它的平方是偶数）。

\vspace{0.5em}
(3) 原命题的否定：
存在一个平行四边形不是中心对称图形。
\end{tcolorbox}
\end{frame}

% ========== 1.5.2 练习 第2题 ==========
\begin{frame}[t]{1.5.2 练习 第2题：写存在量词命题的否定}
\textbf{2.} 写出下列命题的否定：

\vspace{0.5em}
(1) 有些三角形是直角三角形；

\vspace{0.3em}
(2) 有些梯形是等腰梯形；

\vspace{0.3em}
(3) 存在一个实数，它的绝对值不是正数。
\end{frame}

\begin{frame}[t]{1.5.2 练习 第2题 解答}
\begin{tcolorbox}[colback=orange!10,colframe=orange,title=\textbf{解}]
\textbf{规则：} 存在量词命题 $\exists x \in M,\ p(x)$ 的否定为 $\forall x \in M,\ \neg p(x)$。

\vspace{0.8em}
(1) 原命题的否定：
所有三角形都不是直角三角形。

\vspace{0.5em}
(2) 原命题的否定：
所有梯形都不是等腰梯形。

\vspace{0.5em}
(3) 原命题的否定：
任意一个实数，它的绝对值都是正数。

\vspace{0.3em}
\small（注意：实数 $0$ 的绝对值是 $0$，不是正数，所以原命题为真，其否定为假。）
\end{tcolorbox}
\end{frame}

\begin{frame}[c]
  \begin{center}
    {\Huge\calligra Thanks!}
  \end{center}
\end{frame}

\end{document}
